\documentclass[10pt,a4paper]{amsart}
\usepackage[margin=19mm]{geometry}
\usepackage{amsmath,amssymb,mathtools}
\usepackage[scr=rsfs,cal=euler]{mathalfa}
\usepackage{xfrac}
\usepackage[unicode,hidelinks,bookmarks=false]{hyperref}
\newcommand{\ie}{i.e.\ }
\newcommand{\cf}{cf.\ }
\newcommand{\resp}{respectively\ }
\newcommand{\Subsection}{\subsection}
\providecommand{\nobreakdash}{\nobreak\hbox{-}}
\setlength{\emergencystretch}{2em}
\multlinegap0pt
\usepackage{bm}
\usepackage{bbm}
\usepackage{dsfont}
\usepackage{xspace}
\usepackage{cancel}
\usepackage{mathtools}
\usepackage{amsthm}
\makeatletter
\@ifundefined{proposition}{\newtheorem{proposition}{Proposition}}{}
\@ifundefined{theorem}{\newtheorem{theorem}{Theorem}}{}
\makeatother
\title[Relaxation methods for pessimistic bilevel optimization]{Relaxation methods for pessimistic bilevel optimization}
\author{Imane Benchouk, Lateef Jolaoso, Khadra Nachi, Alain Zemkoho}
\date{Source: arXiv:2412.11416v2 (CC BY 4.0)}
\begin{document}
\maketitle

\noindent\emph{Source and licence.} Relaxation methods for pessimistic bilevel optimization. Version arXiv:2412.11416v2 (CC BY 4.0), distributed under the Creative Commons Attribution 4.0 International licence, as recorded in the arXiv licence metadata for this exact version and shown on the abstract page https://arxiv.org/abs/2412.11416v2. Full rights notice: licenses/arxiv-2412.11416-CC-BY-4.0.txt.\par\smallskip

\noindent\textit{Editorial scope.} Counted unit: the Proposition whose source label is \texttt{cs} in the article (source0.tex lines 1168--1172, source label \texttt{cs}), delivered together with the article's own complete original proof of it (source0.tex lines 1173--1224, one \texttt{proof} environment of 52 lines). No mathematics is written, repaired or re-derived: statement and proof are the article's own text line for line. Required context retained uncounted: fprop (lines 885--893); e1 (lines 909--912). Every definition, display and label of those lines is retained unchanged. Omitted from this package and not redistributed: 1--884, 894--908, 913--1167, 1225--4118. Nothing inside a retained excerpt is omitted or altered: every retained body is delivered exactly as the article's lines are written, so no omission receipt of any kind accompanies this package. Citations: the retained excerpts carry no citation command at all, so no bibliography entry of the article is transcribed and no reference is redistributed. Reference adaptation: none was needed and none was made; every reference inside the delivered unit resolves inside it, so no unresolved reference or citation is produced. Printed slips kept literal: no printed word, symbol, hypothesis or number of the counted statement or of its proof was altered. Layout and class: the article's own class is \texttt{svjour3} with packages including fix-cm, csquotes, babel, amsmath, amssymb, inputenc, lscape, algorithm, algorithmic, comment, booktabs, multirow, subfigure, hyperref; the wrapper is the lane's pinned \texttt{amsart} editorial wrapper at 10pt on A4, which restates the theorem-like environments the retained text needs and adds the article's own macro definitions that the retained text needs (none), with extra wrapper packages (bm, bbm, dsfont, xspace, cancel, mathtools). The delivered numbering is the unit's own. Assets: no retained span contains a figure environment, a graphics-inclusion command, a PDF-page inclusion, a table or a TikZ picture; no image file is copied into this package, the package declares no asset dependency and no image file is redistributed with it. Licence: version arXiv:2412.11416v2 of this article is distributed under the Creative Commons Attribution 4.0 International licence, as recorded in the arXiv licence metadata for this exact version and shown on the abstract page https://arxiv.org/abs/2412.11416v2. A full rights notice is delivered at licenses/arxiv-2412.11416-CC-BY-4.0.txt. Characters: every retained excerpt is pure ASCII, so the delivered engine, XeLaTeX with its default Latin Modern text font, reports no missing character.
\begin{theorem} \label{fprop}
Let $\mathcal{R}\in \left\{\mbox{S}, \mbox{LF}, \mbox{SU}, \mbox{KS}\right\}$ then 
\begin{enumerate}
\item[$(a)$]  For all $x\in \mathbb{R}^n$ and $t>0$, $\psi_p(x) \leq \psi^t_{\mathcal{R}}(x)$;
\item[$(b)$] If the function  $t\mapsto\psi_{\mathcal{R}}^{t}(x)$ is  upper semicontinuous at $0^+$ for any $x\in \mathbb{R}^n$ and
	 $(t_{k})\downarrow0$, then for any $x\in\mathbb{R}^{n}$, we have
	$\psi_{\mathcal{R}}^{t_{k}}(x)\rightarrow\psi_{p}(x)$ as $k\rightarrow\infty$.
\end{enumerate}
\end{theorem}

\begin{equation}
F(x,y):=x+y,\text{ \ }X:=[0,1],\text{ \ }Y(x):=\left]  -\infty,1\right]
,\text{ and \ \ }f(x,y):=-xy. \label{e1}
\end{equation}

\begin{proposition} \label{cs}
Let the sequences $(t_{k})\downarrow0$ and $(x^{k}) \rightarrow \bar{x}$ be such that $\mathcal{D}(\bar{x})$ is
nonempty and compact. Let the set-valued mapping $(t,x) \rightrightarrows \mathcal{D}_\mathcal{R}^{t}(x)$ be  upper semicontinuous at $(0^+,\bar{x})$ and $x \rightrightarrows \mathcal{D}(x)$  be lower semicontinuous at $\bar{x}$. Then the property \eqref{e1} holds for any $\mathcal{R}\in \left\{\mbox{S}, \mbox{LF}, \mbox{SU}, \mbox{KS}\right\}$ and it is satisfied for $\mathcal{R} = \text{KDB}$ if in addition the following inequality holds: $\underset{k\rightarrow\infty}{\text{ }\lim\inf}\psi_{p}(x^{k})\leq \underset
{k\rightarrow\infty}{\text{ }\lim\inf}\psi_{\text{KDB}}^{t_{k}}(x^{k})$.
\end{proposition}

\begin{proof}
Assume that for a given $\mathcal{R}\in  \left\{\text{S}, \, \text{LF}, \,\text{KDB}, \, \text{SU},\, \text{KS}\right\}$, the property (\ref{e1}) does not hold. Thus,  there exist $\delta >0$ and a sequence $(y^{k},u^{k})_{k}$ such that $(y^{k},u^{k})\in\mathcal{S}_\mathcal{R}^{t_{k}}(x^{k})$ and
\begin{equation}\label{not}
d\left((y^{k},u^{k}),\;\mathcal{S}_p(\bar{x})\right) \geq \delta \;\mbox{ for all }\; k.
\end{equation}
Therefore, since $(y^{k},u^{k})\in \mathcal{D}_\mathcal{R}^{t_{k}}(x^{k})$, 
$
d\left((y^{k},u^{k}), \,\mathcal{D}(\bar{x})\right)\leq e\left(\mathcal{D}_\mathcal{R}^{t_{k}}(x^{k}
), \,\mathcal{D}(\bar{x})\right).
$
By the upper semicontinuity of the set-valued mapping $(t,x) \rightrightarrows \mathcal{D}_\mathcal{R}^{t}(x)$ at the point $\left(0^+,\bar{x}\right)$,  we have  
\[
\underset{k \rightarrow
		\infty}{\lim }e\left(\mathcal{D}^{t_{k}}_\mathcal{R}(x^{k}), \;\mathcal{D}(\bar{x})\right)=0.
\]
Hence, one can pick a sequence  $(z^{k},w^{k})_k$ in $\mathcal{D}(\bar{x})$ such that
\[
    \underset{k \rightarrow
		\infty}{\lim }\left\Vert y^{k}-z^{k}\right\Vert = 0 \;\, \text{ and } \;\,\underset{k \rightarrow
		\infty}{\lim }\left\Vert
	u^{k}-w^{k}\right\Vert = 0
\]
%\[
%\left\Vert y^{k}-z^{k}\right\Vert \longrightarrow0 \;\, \text{ and }\;\, \left\Vert
%u^{k}-w^{k}\right\Vert \longrightarrow 0
%\]
hold. Due to the compactness of $\mathcal{D}(\bar{x})$, this sequence (up to a subsequence) converges to a point $(\bar{y},\bar
{u})\in\mathcal{D}(\bar{x})$ and so does the (sub)sequence $(y^{k},u^{k})$.

Since $F(\cdot,\cdot)$ is continuous and $\mathcal{D}(\cdot)$ is lower
semicontinuous, the function $\psi_{p}(\cdot)$ is lower semicontinuous at
$\bar{x}$. Let us now prove that $(\bar{y},\bar{u})\in\mathcal{S}_{p}(\bar{x})$ for $\mathcal{R}\in  \left\{\mbox{S}, \, \mbox{LF}, \, \mbox{SU},\, \mbox{KS}\right\}$. Indeed,
from assertion (a) in Theorem \ref{fprop}, 
\begin{equation} \label{psiR}
    F(x^{k},y^{k})-\psi_{p}(x^{k})\geq F(x^{k},y^{k})-\psi_\mathcal{R}^{t_{k}}
(x^{k})\geq0
\end{equation}
as $(y^{k},u^{k})\in\mathcal{S}_\mathcal{R}^{t_{k}}(x^{k})$.  Hence,
\[
F(\bar{x},\bar{y})-\psi_{p}(\bar{x})\geq F(\bar{x},\bar{y})-\underset
{k\rightarrow\infty}{\text{ }\lim\inf}\psi_{p}(x^{k})\geq \underset
{k\rightarrow\infty}{\text{ }\lim\sup}(F(x^{k},y^{k})-\psi_\mathcal{R}^{t_{k}}(x^{k}
))\geq0.
\]
Consequently, $(\bar{y},\bar{u})\in\mathcal{S}_{p}(\bar{x})$ and from the inequality
\[
d\left((y^{k},u^{k}), \;\mathcal{S}_{p}(\bar{x})\right)\leq\left\Vert (y^{k},u^{k})-(\bar
{y},\bar{u})\right\Vert,
\]
$\underset{k\longrightarrow\infty}{\lim}d\left(\left(y^{k},u^{k}\right), \;\mathcal{S}_{p}
(\bar{x})\right)=0$, which is a contradiction to (\ref{not}).  The  poof  is similar for the KDB-relaxation using the additional assumption since the first inequality  \eqref{psiR} fails in this case. \hfill \qed
\end{proof}

\end{document}
